% =============================================================================
%  RESUMO
% =============================================================================
\newpage
\begin{resumo}[Resumo]
\SingleSpacing

Esta obra apresenta uma investigação sistemática e aprofundada sobre a aplicação de gêmeos digitais (\textit{Digital Twins}) na odontologia, com ênfase no desenvolvimento do framework SUS-Twin para simulação biomecânica periodontal no âmbito do Sistema Único de Saúde brasileiro. O livro está estruturado em quatro partes complementares: (i) fundamentos dos gêmeos digitais, abrangendo definições, arquiteturas e aplicações na medicina; (ii) odontologia digital, explorando ortodontia, implantodontia, endodontia, prótese e educação odontológica baseada em simulação; (iii) ecossistemas de inteligência artificial e plataformas abertas, detalhando o OpenCode Ecosystem, seus pipelines de reconstrução 3D, métodos de K-Fold para predição periodontal e gateways IoT para telemetria; e (iv) impacto social e perspectivas futuras, analisando equidade no SUS, ética, privacidade, marcos regulatórios e desafios biomecânicos.

A metodologia empregada combina revisão sistemática da literatura (20+ publicações 2023-2026), modelagem matemática viscoelástica do ligamento periodontal (séries de Prony), validação cruzada K-Fold com RMSE $< 0.15$ mm, e provas de conhecimento zero (ZKP) para auditoria criptográfica de prontuários. O framework SUS-Twin foi implementado em Python com suíte de 11 testes TDD (100\% de aprovação) e dashboard interativo em Streamlit. A análise de Retorno Social do Investimento (SROI) indica um índice médio de $3,5:1$, demonstrando a viabilidade econômico-social da tecnologia para o sistema público de saúde brasileiro.

\textbf{Palavras-chave}: Gêmeos Digitais. Odontologia Digital. Inteligência Artificial. SUS-Twin. OpenCode Ecosystem. Biomecânica Periodontal. Equidade em Saúde.
\end{resumo}
